% Root draft; independent literal audit requested before input.
\subsection{One-sided minimal rows and small categorical acts}

\begin{lemma}[At most two states at each object]
\label{C:two-states-objectwise-act-cups}
Let $S$ be a finite category and let
$F:S^{\mathrm{op}}\to\mathrm{Set}$ be a right act such that
$1\leq |F(c)|\leq2$ at every object. Adjoin one trivial sink
with outgoing hom-sets $F(c)$ from $S$, and call the resulting
act cone $X$. Every product of two positive-degree rational
cohomology classes on $BX$ is zero.
\end{lemma}
\begin{proof}
For every object $c$ take the convex simplex $I(c)$ on the
one- or two-element set $F(c)$, and extend each act map affinely.
These extensions compose strictly, since affine maps on a simplex
are determined by its vertices. Every $I(c)$ is contractible.
The pointed quotient $I(c)/F(c)$ is a point when $|F(c)|=1$
and a circle when $|F(c)|=2$. Every induced map between these
quotients is a pointed constant map, an identity, or a reflection,
after choosing an orientation for each circle.

The degree-two map on every such circle, and the identity of
each point, form a strict natural self-map of this pointed
diagram. Indeed, $t\mapsto2t$ on $\mathbb R/\mathbb Z$
commutes with $t\mapsto-t$, with identities and with pointed
constant maps. The pointed homotopy colimit of $I/F$ is the
cofiber of $\operatorname{hocolim}F\to\operatorname{hocolim}I$.
One can see this by first taking the ordinary homotopy colimit
of the objectwise quotient pushouts and then collapsing its
basepoint section $BS$. Since
$\operatorname{hocolim}I\to BS$ is an equivalence, this cofiber
is the classifying space of the stated act cone.

The reduced cellular chain complex of each quotient is
concentrated in degree one. The natural self-map therefore acts
as multiplication by two on every chain group of the reduced
pointed bar, and hence on all positive rational cohomology.
For positive-degree classes $x,y$, multiplicativity gives
$2xy=4xy$, so $xy=0$.
\end{proof}

\begin{lemma}[Collapsing a source block with one minimal row]
\label{C:one-row-source-block-collapse}
Let $C$ be a finite category partitioned into a full strongly
connected source block $S$ and its full complement $D$.
Assume the minimal ideal of $S$ has trivial maximal group and
exactly one minimal row at every object. No restriction is
imposed on its column counts. Replacing $S$ by one trivial
object and each incident right $S$-act by its orbit set gives
an actual quotient functor $Q:C\to\bar C$ inducing a homotopy
equivalence. Dually a sink block with trivial minimal group
and one minimal column at each object can be collapsed.
\end{lemma}
\begin{proof}
Choose an object $a\in S$ and a minimal idempotent
$e\in S(a,a)$. The minimal Rees description says that
$S(a,i)e$ has exactly one element for every $i$; call it
$u_i:a\to i$. At $i=a$ this element is $e$, and for every
$f:i\to j$ one has
\[
 u_i e=u_i,\qquad f u_i=u_j.
\]
The arrows $u_i$ give a natural transformation from the constant
functor at $a$ to the identity of $S$, proving that $BS$ is
contractible. A component of a natural transformation need not
be an identity, so $u_a=e$ causes no problem.

For a finite right $S$-act $F$, let its element category have
arrows $(i,x)\to(j,y)$ represented by $f:i\to j$ with $x=yf$.
Put $T=F(a)e$ and define
\[
 p(i,x)=x u_i\in T,\qquad s(t)=(a,t).
\]
The equation $f u_i=u_j$ proves that $p$ is constant on every
arrow, while $ps=\mathrm{id}_T$ because $te=t$.
The arrows $u_i$ give a natural transformation
$sp\Rightarrow\mathrm{id}$: its component at $(i,x)$ has
source $(a,xu_i)$ and target $(i,x)$, and its naturality is
exactly the same equation. Thus the element category is
homotopy equivalent to the discrete set $T$, which is its
orbit set. These functors and natural transformations commute
strictly with every morphism of right acts.

For each object $d\in D$ apply this construction to
$F_d(i)=C(i,d)$. Postcomposition by arrows of $D$ preserves
the whole construction. The two-block correspondence homotopy
pushout therefore compares $BC$ with $B\bar C$ by equivalences
on its source block and cross-arrow space, and by the identity
on $BD$. Its homotopy invariance proves the assertion for the
actual orbit-collapse functor. Since no arrow goes from $D$
back into $S$, quotient composition is exactly the induced
$D$-action on these orbit sets. Opposition gives the dual.
\end{proof}

\begin{proposition}\label{C:proper-strong-pair-at-least-seventeen}
Let $C$ be a connected finite category with at most twenty-one
arrows and a proper strongly connected component $S$ on two
objects. If $S$ has at least seventeen arrows, then
$BC\not\simeq\mathbb CP^2$.
Consequently a minimal-arrow projective-plane candidate through
twenty-one arrows has every proper two-object strong component
of size at most sixteen; moreover an extreme such component
with trivial minimal group cannot have one minimal row per
object at a source, or one minimal column per object at a sink.
\end{proposition}
\begin{proof}
Let $\ell$ be the number of objects outside $S$, so $\ell\geq1$.
Contracting $S$ in the connected undirected support leaves a
graph on $\ell+1$ vertices, which needs at least $\ell$ edges.
At least one such edge is incident with $S$. Strong connectivity
forces its incidence to include an arrow at both objects of $S$,
giving at least one additional arrow beyond this support count.
Together with the $\ell$ outside identities,
\[
 |C|\geq |S|+2\ell+1.
\]
For $|S|\geq17$ the bound twenty-one implies $\ell=1$.
Write $M$ for the remaining corner monoid and $t$ for the
number of cross arrows. All these arrows have one direction,
and $t\geq2$. If $t=2$, there is one arrow at each object of
$S$ and the category is the categorical join of its two blocks.
Its positive rational products vanish.

Otherwise $t\geq3$, and the budget forces
$|S|=17$, $|M|=1$, $t=3$. Take opposites if needed to make
$S$ the source. Its two outgoing stalks have sizes one and two,
so Lemma~\ref{C:two-states-objectwise-act-cups} excludes the
nonzero square of $\mathbb CP^2$.

The asserted size bound follows. The final extremal assertion
follows from Lemma~\ref{C:one-row-source-block-collapse}:
such a collapse preserves the homotopy type and strictly
decreases the arrow count, contradicting minimality. As always,
the minimality conclusion is about a category with the fewest
arrows, and does not assert that a nonminimal model cannot
contain the stated component.
\end{proof}
